\paragraph{試験族の二つの $j$ 不変量.}
Hermite--Pfaffian恒等式により $E_Q$ の基底不変性はnetの幾何から分かる。
上の試験族に対応する判別式曲線は
\[
 \Delta_t:\quad\lambda_0^3+4\lambda_1^3+t\lambda_2^3
 -(t+4)\lambda_0\lambda_1\lambda_2=0.
\]
$t=1$ では $E_t,\Delta_t$ ともに滑らかで、$\det G=-3\ne0$ である。
通常の $j(i)=1728$ の正規化では
\begin{align*}
 j(E_t)&=-\frac{(t-8)^3\bigl((t-8)^3-432t^2\bigr)^3}
 {2t^2(t-2)^3(t+16)^6},\\
 j(\Delta_t)&=\frac{(t+4)^3\bigl((t+4)^3+864t\bigr)^3}
 {4t(t-2)^6(t+16)^3}.
\end{align*}
Hesse形 $X^3+Y^3+Z^3+3hXYZ=0$ の公式
$27h^3(8-h^3)^3/(1+h^3)^3$ から得られる
\cite[\S2]{ArtebaniDolgachevHesse}。
いずれも非定数であり、$\mathcal R_3\dashrightarrow Y_0(2)$ の支配性を確認できる。
十八次元の行列空間上の有理関数の非定数性なので、元の四次Gram曲線が滑らかな
稠密開部分へ制限しても変わらない。
二つの $j$ が満たす次数二のモジュラー方程式は恒等式として再検査できるが、
被覆の存在とtrace零線束の同定は第\ref{subsec:notes-hermite-v27}節の幾何学的証明に依拠する。
